% Preámbulo: paquetes y estilos de la memoria. Compila con pdflatex.

% --- Codificación, idioma y tipografía ---
\usepackage[utf8]{inputenc}
\usepackage[T1]{fontenc}
\usepackage{lmodern}
\usepackage{textcomp}
% es-noshorthands: las comillas " y los signos < > no se vuelven activos
\usepackage[spanish,es-tabla,es-noshorthands]{babel}
\usepackage[a4paper,margin=2.5cm,headheight=14pt]{geometry}
\usepackage{placeins}
\usepackage{microtype}

% Caracteres UTF-8 del texto que inputenc no conoce por defecto
\DeclareUnicodeCharacter{2192}{\ensuremath{\rightarrow}}
\DeclareUnicodeCharacter{2212}{\ensuremath{-}}
\DeclareUnicodeCharacter{2248}{\ensuremath{\approx}}

% --- Color, tablas, listas, figuras ---
\usepackage{xcolor}
% Colores de marca de Unicaja
% (sacados de la web pública de Unicaja, no de un manual de marca).

% Papeles principales
\definecolor{unicajaPetroleo}{HTML}{005265}        % cabecera, saldo, botones
\definecolor{unicajaPetroleoOscuro}{HTML}{00323A}  % degradados, texto sobre lima
\definecolor{unicajaVerde}{HTML}{278600}           % acento, estado activo
\definecolor{unicajaLima}{HTML}{95E891}            % sello u, botones secundarios
\definecolor{unicajaAmbar}{HTML}{EDAF21}           % segundo color de los gráficos
\definecolor{unicajaTinta}{HTML}{14232B}           % texto principal

% Estados (éxito, información, peligro, aviso) y sus tintas claras
\definecolor{unicajaExito}{HTML}{048848}
\definecolor{unicajaInfo}{HTML}{026169}
\definecolor{unicajaPeligro}{HTML}{CC2830}
\definecolor{unicajaAviso}{HTML}{825300}
\definecolor{unicajaExitoClaro}{HTML}{E6F3ED}
\definecolor{unicajaInfoClaro}{HTML}{E6F3F4}
\definecolor{unicajaPeligroClaro}{HTML}{FAEAEA}
\definecolor{unicajaAvisoClaro}{HTML}{FDF5E6}

% Neutros: fondo y bordes
\definecolor{unicajaFondo}{HTML}{F6F9F9}
\definecolor{unicajaBordeClaro}{HTML}{E5EEF0}
\definecolor{unicajaBorde}{HTML}{D0DAD6}

\usepackage{graphicx}
\usepackage{array}
\usepackage{booktabs}
\usepackage{longtable}
\usepackage{tabularx}
\usepackage{enumitem}
\usepackage{caption}
\captionsetup{font=small,labelfont={bf,color=unicajaPetroleo}}
\usepackage{tikz}
\usetikzlibrary{positioning,fit,backgrounds,arrows.meta,calc,shapes.geometric}
\usepackage[most]{tcolorbox}
\usepackage{listings}
\usepackage{fancyhdr}
\usepackage{titlesec}
\usepackage{needspace}

% --- Enlaces y rutas ---
\usepackage{url}
\usepackage[colorlinks=true,linkcolor=unicajaPetroleo,urlcolor=unicajaVerde,
            citecolor=unicajaPetroleo,pdfauthor={Florin Emanuel Todor Gliga},
            pdftitle={Unia, el asistente de Unicaja. Memoria técnica}]{hyperref}

% \ruta{...}: rutas e identificadores con guiones bajos, sin escapar y con
% cortes de línea permitidos en / . _ (solo para textos sin espacios ni
% caracteres no ASCII).
\DeclareUrlCommand\ruta{\urlstyle{tt}}
% \cod{...}: código en línea normal (hay que escapar los caracteres especiales).
% Permite estirar un poco el espacio entre caracteres de \ruta para justificar sin desbordar.
\Urlmuskip=0mu plus 2mu
\setlength{\emergencystretch}{2em}
\newcommand{\cod}[1]{\texttt{#1}}

% --- Títulos de apartado con el color de Unicaja ---
\titleformat{\section}{\Large\bfseries\color{unicajaPetroleo}}{\thesection.}{0.6em}{}
\titleformat{\subsection}{\large\bfseries\color{unicajaPetroleoOscuro}}{\thesubsection}{0.6em}{}
\titlespacing*{\section}{0pt}{1.6ex plus .5ex}{1ex}
\titlespacing*{\subsection}{0pt}{1.4ex plus .4ex}{.8ex}
% Sin \needspace en los títulos: dentro de \titleformat dejaba páginas a medias y descolocaba las longtable.

% --- Párrafos ---
\setlength{\parindent}{0pt}
\setlength{\parskip}{0.55em}

% --- Tablas: un poco más de aire y columnas ajustadas ---
\renewcommand{\arraystretch}{1.18}
\newcolumntype{L}{>{\raggedright\arraybackslash}X}
\newcolumntype{P}[1]{>{\raggedright\arraybackslash}p{#1}}
\setlength{\LTpre}{0.4em}
\setlength{\LTpost}{0.4em}

% --- Listas compactas ---
\setlist{itemsep=2pt,topsep=3pt}

% --- Encabezado y pie ---
\pagestyle{fancy}
\fancyhf{}
\fancyhead[L]{\small\color{unicajaPetroleo}Unia · Memoria técnica}
\fancyhead[R]{\small\color{unicajaPetroleo}Reto ``Habla con tu dinero''}
\fancyfoot[C]{\small\thepage}
\renewcommand{\headrulewidth}{0.4pt}
\renewcommand{\headrule}{{\color{unicajaPetroleo}\hrule height \headrulewidth width \headwidth}}

% --- Código: listings con acentos UTF-8 ---
\lstdefinestyle{unia}{
  basicstyle=\ttfamily\small,
  backgroundcolor=\color{unicajaFondo},
  frame=single,
  rulecolor=\color{unicajaBorde},
  framesep=4pt,
  xleftmargin=4pt,
  xrightmargin=4pt,
  keywordstyle=\color{unicajaPetroleo}\bfseries,
  commentstyle=\color{unicajaAviso}\itshape,
  stringstyle=\color{unicajaExito},
  showstringspaces=false,
  columns=fullflexible,
  keepspaces=true,
  breaklines=true,
  upquote=true,
  literate={á}{{\'a}}1 {é}{{\'e}}1 {í}{{\'i}}1 {ó}{{\'o}}1 {ú}{{\'u}}1
           {Á}{{\'A}}1 {É}{{\'E}}1 {Í}{{\'I}}1 {Ó}{{\'O}}1 {Ú}{{\'U}}1
           {ñ}{{\~n}}1 {Ñ}{{\~N}}1 {ü}{{\"u}}1 {…}{{\ldots}}1
}
\lstset{style=unia}

% --- Caja destacada (para la idea central del diseño) ---
\newtcolorbox{destacado}{colback=unicajaInfoClaro,colframe=unicajaPetroleo,
  boxrule=0.6pt,arc=2pt,left=6pt,right=6pt,top=4pt,bottom=4pt}
